\ifthenelse{\equal{\Langue}{FR}}{
\chapter{Recommandations pour les TPs}
}{
\chapter{Recommendations for the Labs}
}

\ifthenelse{\equal{\Langue}{FR}}{
Le but de ces travaux pratiques est d'illustrer le cours de conversion d'énergie de deuxième année par l'étude de quatre manipulations:
}{
The purpose of these labs is to illustrate the second year course on energy conversion through the study of four experiments:
}

\begin{itemize}

\ifthenelse{\equal{\Langue}{FR}}{
\item Hacheur Buck (simulations avec LTSpice)
\item Hacheur 4 quadrants
\item Alimentation à découpage Flyback
\item Alimentation à découpage Forward
}{
\item Buck converter (simulations with LTSpice)
\item Four-quadrant chopper
\item Flyback switching power supply
\item Forward switching power supply
}
\end{itemize}

\ifthenelse{\equal{\Langue}{FR}}{
Afin de préparer ces séances, il est nécessaire de connaître le cours correspondant et d'avoir lu le texte complet du TP afin de prévoir les différentes étapes de l'étude. En particulier, lorsqu'une comparaison avec les courbes théoriques est demandée, il est nécessaire de les apporter en séance et de savoir les adapter au TP\@.

Les thèmes de TP se déroulent en rotation: il est donc impératif de connaître son numéro de binôme avant de venir en TP\@.
}{
In order to prepare for these labs, it is necessary to know the corresponding course and to have read the complete lab text in order to plan the different steps of the study. In particular, when a comparison with theoretical curves is requested, it is necessary to bring them to the lab and to know how to adapt them to the lab.

The lab topics are conducted in rotation: it is therefore imperative to know your binomial number before coming to the lab.
}


\vspace{0.5cm}
\hspace{-1cm}
\begin{tabular}{|c | c | c | c | c | c | c | c|}
\hline
\ifthenelse{\equal{\Langue}{FR}}{
& \textbf{Binôme 1} & \textbf{Binôme 2} & \textbf{Binôme 3} & \textbf{Binôme 4} & \textbf{Binôme 5} & \textbf{Binôme 6} & \textbf{Binôme 7}\\
}{
& \textbf{Pair 1} & \textbf{Pair 2} & \textbf{Pair 3} & \textbf{Pair 4} & \textbf{Pair 5} & \textbf{Pair 6} & \textbf{Pair 7}\\
}
\hline
\ifthenelse{\equal{\Langue}{FR}}{
\textit{Séance 1} & Hacheur 4Q & Hacheur 4Q & Hacheur 4Q & LTSpice & LTSpice & LTSpice & LTSpice\\
}{
\textit{Session 1} & 4Q Chop. & 4Q Chop. & 4Q Chop. & LTSpice & LTSpice & LTSpice & LTSpice\\
}
\hline
\ifthenelse{\equal{\Langue}{FR}}{
\textit{Séance 2} & LTSpice & LTSpice & LTSpice & Hach 4Q & Hach. 4Q & Hach. 4Q & Hach. 4Q\\
}{
\textit{Session 2} & LTSpice & LTSpice & LTSpice & 4Q Chop. & 4Q Chop. & 4Q Chop. & 4Q Chop.\\
}
\hline
\textit{Séance 3} & Forward & Forward & Forward & Flyback & Flyback & Flyback & Flyback\\
\hline
\textit{Séance 4} & Flyback & Flyback & Flyback & Forward & Forward & Forward & Forward\\
\hline
\end{tabular}

\vspace{0.5cm}
\ifthenelse{\equal{\Langue}{FR}}{
Suivant le laboratoire où se déroulent les séances de TP, quelques différences peuvent apparaître sur le câblage des machines lors du TP hacheur quatre quadrants ainsi que sur les grandeurs nominales de ces machines qui figurent sur leur plaque signalétique et qui devront donc être soigneusement notées.
}{
Depending on the laboratory where the lab sessions take place, some differences may appear in the wiring of the machines during the four-quadrant chopper lab as well as in the nominal values of these machines which are shown on their nameplate and which must therefore be carefully noted.
}
\vspace{0.5cm}

\ifthenelse{\equal{\Langue}{FR}}{
\textbf{\underline{A l'issue de chaque séance}, un compte-rendu devra être rendu au professeur. L'évaluation tiendra compte à la fois de la qualité du compte rendu et du comportement lors de la séance elle-même.}
}{
\textbf{\underline{At the end of each session}, a report must be submitted to the teacher. The evaluation will take into account both the quality of the report and the behavior during the session itself.}
}